\documentclass[12pt]{article}

\usepackage[english,russian]{babel}

\begin{document}

\section{Оценка результатов.}

Теория отслеживания --- бурно развивающаяся область теории динамических систем.
Однако достаточно мало известно об обратных свойствах отслеживания, и достаточно мало известно о свойствах отслеживания групповых действий. Так как работа А. В. Фадеева развивает два эти не слишком разработанных направления, она представляет несомненный интерес. В работе рассматриваются действия группы Баумслага-Солитара. Эта группа часто используется как источник различных контрпримеров --- в частности, С. Б. Тихомиров и А. В. Осипов показали, что в задаче прямого отслеживания она занимает некоторое промежуточное положение между нильпотентными группами и свободными группами. Кроме того, С. Ю. Пилюгин доказал вариацию леммы об обратном отслеживании для групповых действий нильпотентных групп. В данной работе А. В. Фадеев построил нетривиальные примеры действий группы Баумслага-Солитара, не обладающие обратным свойством отслеживания. Думаю, что эти результаты представляют большой интерес, потому что подобные результаты об обратном отслеживании групповых действий ранее не публиковались.

\section{Список замечаний.}

У меня есть следующие замечания по поводу работы:

\begin{enumerate}
	\item Мне кажется хорошей идей кратко неформально формулировать результаты других работ, касающихся свойств отслеживания групповых действий, потратив на это не больше абзаца. Так читатель сможет понять, что для каких-то нетривиальных групповых действий удалось установить обратное отслеживание, узнать, что известно про прямое отслеживание для групповых действий группы Баумслага-Солитара. Это поможет читателю лучше оценить результаты данной статьи.
	\item стр. 2, 2-й абзац сверху. Имеет смысл сказать, что в $BS(1, n)$ $n$ --- это натуральное число.
	\item В данной статье в формулах стабильно опускаются кванторы всеобщности. Мне кажется, их лучше ставить.
	\item Стр. 7, 2-й абзац сверху, имеет смысл формально определить $p_1$ и $p_2$.
	\item Складывается ощущение, что для доказательства основных утверждений можно было брать не произвольные точки $p$, а просто $p = 0$ или $p = (0, 0)$ соответственно. Это бы немного упростило формулы.
\item Имеет смысл перепроверить англоязычный абстракт в конце статьи:

Например, в абстракте есть фраза ''the theory of inverse shadowing was developed''.

Кажется, что в этой фразе предполагается, что развитие теории обратного отслеживания уже завершено. Складывается ощущение, что лучше было бы написать:
''the theory of inverse shadowing was being developed''.

Кроме того, кажется, что в фразе
''the Baumslag-Solitar group which is often considered as a source of counterexamples.''
не хватает запятой перед which, так как это скорее non-defining relative clause.	

\section{Вывод.}

Мне кажется, что это очень хорошая и интересная работа. Я считаю, что она достойна того, чтобы быть опубликованной в Вестнике СПбГУ. По крайней мере я читал ее с удовольствием.

\end{enumerate}

\end{document}